% ============================================================
% chapter_03_methods.tex - Chapter 3, Methods
%
% Purpose: describes the LSEQM+DL hybrid bias correction framework,
%          the validation protocol, and the metric catalogue.
% Sections: 3.1 Time, Place, and Tools, 3.2 Study Area and Data,
%           3.3 Four-Stage Hybrid Pipeline (LSEQM+DL: LS, EQM/Gamma, GPD tail, CNN),
%           3.4 Station-Density-Aware Blending,
%           3.5 Validation Framework (BMKG out-of-sample),
%           3.6 Verification Metric Catalogue (four-pillar + CQI),
%           3.7 Reproducible Implementation Design.
% Included by main_thesis.tex.
% License: CC BY 4.0 (see ../LICENSE).
% ============================================================

\chapter{Methods}\label{chap:methods}

This chapter describes the study area and data
(Section~\ref{sec:methods-data}), the four-stage hybrid pipeline
LSEQM+DL, comprising Linear Scaling (LS), Empirical Quantile Mapping (EQM), 
Generalized Pareto Distribution (GPD), and Convolutional Neural Network (CNN) 
refinement that constitutes the Deep Learning (DL) stage (Section~\ref{sec:methods-pipeline}). It then sets out
the station-density-aware blending design
(Section~\ref{sec:methods-blending}), the out-of-sample validation
framework (Section~\ref{sec:methods-validation}), the verification
metric catalogue (Section~\ref{sec:methods-metrics}), and the
reproducible implementation that supports independent verification
(Section~\ref{sec:methods-reproducibility}). The algorithmic details
summarised here are extended from a parallel manuscript
\citep{ref:istanto2026mdpi}\MDPIstatus; the present chapter adds
the reproducibility design and the full 31-metric verification
framework that the parallel manuscript reports in compressed form.


% ------------------------------------------------------------
\section{Time, Place, and Tools}\label{sec:methods-setup}
% ------------------------------------------------------------

This research was conducted at the Climatology Laboratory of the
Department of Geophysics and Meteorology, Faculty of Mathematics and
Natural Sciences, Bogor Agricultural University, from
March 2024 to February 2026.

Data processing and model development used the Python scientific
computing stack; the pinned package versions are listed in
Section~\ref{sec:methods-reproducibility}. Code was written and run in
Visual Studio Code on a local workstation and re-run on Google Colab
for cloud-based reproducibility checks, and the thesis was typeset in
LaTeX. The schematic figures, the conceptual workflow and the Stage~4
architecture diagram, were prepared in Inkscape.


% ------------------------------------------------------------
\section{Study Area and Data}\label{sec:methods-data}
% ------------------------------------------------------------

Three datasets enter the framework. The Late Run variant of the
Integrated Multi-satellitE Retrievals for Global Precipitation
Measurement (IMERG) Version 7 (V07) is the product to be
corrected. The CPC Unified Gauge-Based Analysis of Global Daily
Precipitation (CPC-UNI) is the correction reference. The
Indonesian Meteorological, Climatological, and Geophysical Agency
(BMKG) operates the independent station network used for
out-of-sample validation. Each is described in turn, with explicit
attention to the BMKG filter logic that reduces the 180 archived
stations to the 172 retained for the validation reported in
Chapter~\ref{chap:results}.

\subsection{Study Area}

The Indonesian archipelago spans approximately 47 degrees of
longitude from $95^{\circ}$E (Sabang, northwestern Sumatra) to
$141^{\circ}$E (Merauke, eastern Papua) and 17 degrees of
latitude from $6^{\circ}$N to $11^{\circ}$S. It occupies the
centre of the Maritime Continent, where the convergence of the
Pacific and Indian Ocean basins drives a deep convective regime
and where the seasonal cycle is governed by the
Asian-Australian monsoon and the migration of the Inter-Tropical
Convergence Zone \citep{ref:aldrian2003}. Daily-mean precipitation is
high across the archipelago, but the spatial pattern is uneven:
dense convective rainfall over Sumatra, Java, and Kalimantan
contrasts with the comparatively drier eastern provinces. The
convective diurnal cycle peaks at local 14:00 to 17:00 (UTC+7 to
UTC+9), a detail that matters for daily aggregation because the
IMERG calendar day is defined in UTC and therefore splits the
local convective afternoon across two calendar days. The gauge
network is concentrated in Java and southern Sumatra; coverage is
sparse over Papua, Maluku, and eastern Nusa Tenggara
\citep{ref:chen2008}. This spatial gradient motivates the
density-aware blending design described in
Section~\ref{sec:methods-blending}.

\subsection{IMERG Late Run}

IMERG V07 \citep{ref:hou2014, ref:huffman2020} merges passive
microwave retrievals from the Global Precipitation Measurement
(GPM) constellation with geostationary infrared estimates and
monthly gauge analyses to produce a quasi-global precipitation
product at $0.1^{\circ}$ (approximately 11~km) and 30-minute
resolution. Three latency variants are released: Early
(approximately 4~h), Late (approximately 14~h), and Final
(approximately 3.5~months). The framework operates on the Late
Run daily product (IMERG-L, \texttt{GPM\_3IMERGDL.07};
\citealp{ref:imergdl07}) because its 14-hour latency fits within
the same-day decision window of operational services such as
flash-flood warning and reservoir scheduling. The study period
spans 2001-01-01 to 2025-12-31, providing 25 years of daily
fields aggregated from the half-hourly native inputs. Native
IMERG-L files are downloaded as NetCDF and subset to the
Indonesian bounding box before processing. The same Late Run is
also used at its native half-hourly resolution
(\texttt{GPM\_3IMERGHHL.07}; \citealp{ref:imerghhl07}) for the
window-offset timing diagnostic in
Section~\ref{sec:methods-window}; using the gauge-free Late Run
at both resolutions keeps the diagnostic consistent with the
product the pipeline corrects.

\subsection{CPC Unified Gauge-Based Analysis}

CPC-UNI \citep{ref:chen2008} is a gauge-based daily precipitation
analysis produced by the National Oceanic and Atmospheric
Administration (NOAA) Climate Prediction Center on a
$0.5^{\circ}$ (approximately 55~km) global land grid. The product
interpolates daily station observations using optimal
interpolation with a station-density-aware correlation function.
Indonesian contributions come from BMKG and other Global
Telecommunication System (GTS) reports; the contributing station
count varies day to day with GTS reporting completeness, and the
contributing network is concentrated in Java and Sumatra. The
contributing observations follow the World Meteorological
Organization (WMO) local-day morning-observation convention of
their reporting stations, so the CPC-UNI daily total assigned to
date~$d$ at a given grid cell represents a 24-hour accumulation
ending at the local morning observation on date~$d$ rather than a
UTC-day total \citep{ref:cpcuni_readme}; this convention is
relevant to the Pearson~$r$ ceiling discussed in
Section~\ref{sec:discussion-subdaily}.
CPC-UNI is documented to systematically undercatch heavy
rainfall \citep{ref:chen2008}, reflecting both the smoothing introduced by interpolation
onto a $0.5^{\circ}$ grid and the under-representation of
mountain and coastal sites in the contributing network. Despite
these limitations, CPC-UNI is the most widely used gauge-based
daily reference for satellite product evaluation over the tropics
and serves as the correction target throughout this study. The
same 2001 to 2025 window is used.

\subsection{BMKG Validation Network}

BMKG operates a national synoptic station network whose daily
precipitation archive serves as the independent validation
reference. Independence here has a specific meaning. Some BMKG
stations contribute to the global CPC-UNI analysis, so BMKG
observations are not independent of CPC-UNI in a strict
statistical sense. However, the validation pipeline compares
per-station daily observations against the corrected satellite
product extracted at each station's pixel, bypassing the CPC
interpolation step entirely. Validation therefore tests whether
the corrected satellite reproduces the original point
observations, not whether it reproduces the gridded reference.

The BMKG network reports in near-real-time, but the record
extracted for this study covers 180 stations with daily
observations from 2001-01-01 to 2021-12-31. Missing days are
sentinel-coded as $8888.0$ in the source file; after sentinel
removal the extracted record yields 880{,}702 valid daily
observations. Two filters are applied before validation:

\begin{enumerate}
    \item \textit{Spatial mapping filter.} The station's reported
    latitude and longitude must fall within the Indonesian land
    mask used by the satellite product. A small number of
    stations near the maritime border or with mis-recorded
    coordinates fail this check.
    \item \textit{Minimum paired-day filter.} After sentinel
    removal and pairing against the corrected satellite series at
    the station pixel, the station must have at least 30 paired
    daily observations in any dekadal window used for validation.
    Stations with very sparse records fail.
\end{enumerate}

Together the two filters drop 8 of the 180 archived stations,
leaving 172 with usable data for the per-station verification
reported in Chapter~\ref{chap:results}. The 8 dropped stations
fall into two categories. A small subset of stations report
coordinates that lie outside the Indonesian land mask used by
the satellite product, typically because the station is on a
small offshore islet not covered by the $0.1^{\circ}$ mask grid,
or because the coordinate string in the archive has a recording
error. The remaining dropped stations satisfy the spatial filter
but yield fewer than 30 paired daily observations after sentinel
removal, generally reflecting extended periods of equipment
failure, transmission gaps, or stations that came online only
late in the 2001 to 2021 window. The retained network spans the
seven major regions of the archipelago: Sumatra, Kalimantan,
Sulawesi, Java, Bali and Nusa Tenggara, Maluku, and Papua. The
full archived network is shown in Figure~\ref{fig:station-map},
which illustrates the gauge-density gradient referred to
throughout the remainder of the chapter: the network is
concentrated over Java and southern Sumatra, with sparser
coverage over Papua, Maluku, and eastern Nusa Tenggara.

\begin{figure}[H]
\centering
\includegraphics[width=\figThreeOneWidth\textwidth]{fig_thesis_03_station_map.png}
\caption[Spatial distribution of the 180 archived BMKG stations, color-coded by region]{Spatial distribution of the 180 archived BMKG stations, color-coded by region.\label{fig:station-map}}
\end{figure}


% ------------------------------------------------------------
\section{Four-Stage Hybrid Pipeline (LSEQM+DL)}\label{sec:methods-pipeline}
% ------------------------------------------------------------

The framework is built on a one-bias-dimension-per-stage design.
Section~\ref{sec:methods-data} described three bias modes that
IMERG-L carries against the gauge reference: a mean offset, a
distributional mismatch in the bulk, and a truncated upper tail.
Each of these has a matched correction stage: LS for the mean,
EQM with a Gamma fit for the bulk distribution, and GPD tail
adjustment for the extremes. A fourth stage, a CNN, addresses a
limitation that none of the marginal stages can reach by
themselves: the lack of spatial coherence in the extreme pixels
of a per-pixel statistical correction. The four stages execute
in sequence; the output of one stage becomes the input of the
next. All fitting is performed per pixel and per dekadal window,
with the 36 dekadal windows pooled across 2001 to 2025 giving
approximately 250 paired daily observations per pixel
per window. Figure~\ref{fig:pipeline} shows the conceptual
workflow, in which the final stage blends the CNN refinement with
the LSEQM result using a station-density-aware weight
(Section~\ref{sec:methods-blending}).

% Fig 3.2: dedicated page, image rotated 90 deg to occupy the page height
% as its new width (same approach as Fig 3.3 below). The PNG itself is the
% original landscape diagram; LaTeX handles the rotation.
\begin{figure}[p]
\centering
\rotatebox{90}{%
  \includegraphics[width=\figThreeTwoWidth\textheight]{fig_thesis_03_pipeline.png}%
}
\caption[Conceptual workflow of the LSEQM+DL hybrid bias correction framework]{Conceptual workflow of the LSEQM+DL hybrid bias correction framework.\label{fig:pipeline}}
\end{figure}


\subsection{Mean Bias Correction: Linear Scaling}\label{sec:methods-ls}

The first stage addresses the mean offset between the satellite
and the gauge reference. For each CPC 0.5$^{\circ}$ cell~$c$ and
each dekadal window, a multiplicative correction factor is
computed from the pooled means of the satellite and reference
samples \citep{ref:teutschbein2012}:
\begin{equation}
P_{\text{LS}}(x, y, t) =
P_{\text{IMERG}}(x, y, t)
\cdot \frac{\overline{P}_{\text{CPC}}(c)}{\overline{P}_{\text{IMERG}}(c)},
\label{eq:ls}
\end{equation}
where $\overline{P}_{\text{CPC}}$ and $\overline{P}_{\text{IMERG}}$
are the pooled-dekad means of the gauge and satellite precipitation
within cell~$c$, respectively. The factor is applied to every
0.1$^{\circ}$ pixel that falls inside cell~$c$, following the
Bias Correction Spatial Disaggregation (BCSD) strategy of
\citep{ref:wood2004}. LS corrects the dekadal mean by construction;
it does not change the rank ordering of values within a pixel,
nor does it touch the distribution shape or the tail. Subsequent
stages address these remaining bias dimensions on the LS-corrected
series.


\subsection{Distribution Alignment: Empirical Quantile Mapping with Gamma}\label{sec:methods-eqm}

The second stage aligns the full marginal distribution of the
LS-corrected satellite series to the gauge reference. A
two-parameter Gamma distribution is fitted by maximum likelihood
to the wet-day subset of each pixel and each dekadal window,
yielding Gamma cumulative distribution functions (CDFs)
$F_{\text{IMERG},\Gamma}$ and $F_{\text{CPC},\Gamma}$
\citep{ref:gudmundsson2012, ref:themessl2012}. The corrected
value at day~$t$ is then obtained by quantile composition:
\begin{equation}
P_{\text{EQM}}(x, y, t) =
F_{\text{CPC},\Gamma}^{-1}\!
\left(F_{\text{IMERG},\Gamma}\!\left(P_{\text{LS}}(x, y, t)\right)\right).
\label{eq:eqm}
\end{equation}
The Gamma family is the standard parametric model for
daily precipitation on wet days. Its two parameters describe the
shape and the scale of the bulk distribution, where most of the
sample mass concentrates.

Because the gauge reference is coarser than the satellite product,
the CPC parameters cannot be applied at the IMERG resolution
directly. CPC-UNI is native to a $0.5^{\circ}$ grid while IMERG-L is
$0.1^{\circ}$, so each CPC cell spans a $5 \times 5$ block of IMERG
pixels. Following the bias-correction and spatial-disaggregation
principle \citep{ref:wood2004}, the Gamma, GPD, and dry-day
parameters are fit once at the native $0.5^{\circ}$ resolution and
bilinearly interpolated to the $0.1^{\circ}$ grid, rather than
regridding the gauge field itself; each $0.1^{\circ}$ pixel is then
quantile-mapped against its own interpolated reference distribution.
This disaggregates the correction, not the data: the satellite
retains its $0.1^{\circ}$ spatial detail and day-to-day sequence,
while the bilinear interpolation of the parameters avoids the
block-boundary artefacts that nearest-neighbour replication of a
$0.5^{\circ}$ field would impose. Figure~\ref{fig:bcsd} illustrates
the disaggregation on the Bali subdomain: panel~(a) shows that one
$0.5^{\circ}$ CPC cell covers a $5 \times 5$ block of $0.1^{\circ}$
pixels, panel~(b) the wet-day-mean parameter field fit at the native
$0.5^{\circ}$ resolution, panel~(c) the same field after bilinear
interpolation to $0.1^{\circ}$, and panel~(d) the per-pixel quantile
mapping of one $0.1^{\circ}$ pixel against its interpolated CPC
distribution.

\begin{figure}[H]
\centering
\includegraphics[width=\textwidth]{fig_thesis_03_bcsd_disaggregation.png}
\caption[Bias-correction spatial disaggregation of the CPC reference parameters from 0.5 to 0.1 degree]{Bias-correction spatial disaggregation of the CPC reference parameters from $0.5^{\circ}$ to $0.1^{\circ}$ on the Bali subdomain.\label{fig:bcsd}}
\end{figure}

Dry-day handling uses wet-day-frequency matching
\citep{ref:cannon2015}: the corrected series is constrained to
reproduce the gauge wet-day fraction within each dekadal window.
Days with $P_{\text{LS}} = 0$ are kept dry; days with non-zero
LS values that fall below the gauge wet-day threshold are
reclassified to dry. This mechanism is the source of the
event-detection trade-off discussed in
Chapter~\ref{chap:results}: matching the dry-day frequency
suppresses spurious light-rain events at the cost of reduced
aggregate probability of detection, but with the benefit of a
reduced false alarm ratio.

The Gamma family fits the bulk of the wet-day sample by design,
which means its parameters are dominated by light and moderate
rainfall, the segment that contains most of the sample mass. The
upper tail, the segment that is operationally meaningful for
flood and drought thresholds, is fit only as a by-product and is
systematically too short. The next stage corrects this.


\subsection{Extreme Tail Preservation: Generalized Pareto Distribution}\label{sec:methods-gpd}

\medskip\noindent\textbf{Motivation.}
The Gamma fit in Stage~2 is a two-parameter distribution fitted
to the full wet-day sample. By construction its parameters are
dominated by the bulk of the sample, that is, the light to
moderate rainfall that contains most of the points. The upper
tail, which is the segment operationally meaningful for flood
and drought thresholds, is fit only as a by-product of forcing
the Gamma curve through the bulk, and for tropical precipitation
is systematically too short. Gamma cannot represent the heavy
tail that real tropical convective rainfall exhibits. Stage~3
grafts a GPD onto the upper
tail, fitted explicitly to exceedances above a high threshold.
The GPD is the canonical extreme-value model in the
peaks-over-threshold framework: under mild regularity conditions
the distribution of exceedances above a high threshold converges
to a GPD with shape and scale parameters $(\xi, \sigma)$
\citep{ref:coles2001}, with cumulative distribution function
\begin{equation}
F_{\text{GPD}}(x; u, \sigma, \xi) =
1 - \left[1 + \xi\,\frac{x - u}{\sigma}\right]^{-1/\xi},
\quad x > u,
\label{eq:gpd}
\end{equation}
where $u$ is the threshold. The shape parameter $\xi$ controls
tail heaviness: $\xi > 0$ indicates a heavy Pareto-like tail
typical of tropical convective extremes, $\xi = 0$ an exponential
tail, and $\xi < 0$ a bounded tail.

\medskip\noindent\textbf{Method.}
For each pixel and each dekadal window, the GPD fit proceeds in
seven steps.

\begin{enumerate}
    \item Extract the wet-day subset of the 25-year sample
    (approximately 250 paired days multiplied by the wet-day fraction).
    \item Reuse the Gamma fit from Stage~2 as the bulk model.
    \item Compute the 80th-percentile threshold of the wet
    sample.
    \item Extract exceedances (wet observations above the
    threshold), typically 50 points if the wet sample has 250
    entries (250 times 0.20).
    \item Fit GPD parameters $(\xi, \sigma)$ to the exceedances
    by maximum likelihood with the location parameter fixed at
    the threshold \citep{ref:coles2001}; the parameters are
    averaged across five subsamples to reduce variance on the
    small per-pixel exceedance set (see
    Section~\ref{sec:methods-gpd} subsection on subsample
    aggregation below).
    \item Above the 80th percentile, replace the Gamma-based
    quantile from Stage~2 with the GPD-based quantile.
    \item Cap the final value at the 99.9th percentile of the
    CPC reference distribution as a safety net against GPD
    over-extrapolation in the extreme tail.
\end{enumerate}

The replacement in step~6 is a quantile remap through the
reference tail. For a satellite value $P_{\text{LS}}$ above the
satellite threshold $u_{\text{sat}}$, the conditional exceedance
probability within the satellite tail is
\begin{equation}
p_c = \frac{F_{\text{sat}}(P_{\text{LS}}) - F_{\text{sat}}(u_{\text{sat}})}
           {1 - F_{\text{sat}}(u_{\text{sat}})},
\label{eq:gpd-pc}
\end{equation}
where $F_{\text{sat}}$ is the satellite wet-day cumulative
distribution. This probability is then mapped onto the reference
tail through the inverse of the reference GPD cumulative
distribution $G_{\text{ref}}$, anchored at the reference
threshold $u_{\text{ref}}$:
\begin{equation}
P_{\text{LSEQM}} = u_{\text{ref}} + G_{\text{ref}}^{-1}(p_c).
\label{eq:gpd-map}
\end{equation}
In the implementation the reference GPD is fitted to the
threshold exceedances $x - u_{\text{ref}}$ with its location
pinned to zero, so $G_{\text{ref}}^{-1}(p_c)$ returns an
exceedance above $u_{\text{ref}}$ and the threshold is added
back to recover the absolute value.
Equations~\eqref{eq:gpd-pc} and~\eqref{eq:gpd-map} preserve the
rank of the satellite value within its own tail and place it at
the matching quantile of the reference tail, so the heavy-tail
shape of the corrected field follows the reference GPD rather
than the truncated Gamma extrapolation of Stage~2.

\medskip\noindent\textbf{Threshold choice and sample-size considerations.}
The 80th-percentile threshold is unconventionally low for a GPD
fit. Standard practice in extreme-value analysis uses the 90th
to 99th wet-day percentile, justified per region by mean
residual life plots or parameter-stability diagnostics
\citep{ref:coles2001}. The present study made an empirical
trade-off in the other direction: a 90th-percentile choice would
yield only about 25 exceedances per pixel-dekad at the typical
wet-day fraction, on the edge of the sample size at which the
GPD Maximum Likelihood Estimation (MLE) becomes operationally noisy \citep{ref:hoskingwallis1987};
the 80th-percentile choice yields about 50 exceedances and
trades cleanliness of the extreme regime for sample-size
stability. Lower thresholds (60th or 70th) would yield still
more exceedances but conflate the extreme tail with moderate
rain that the Gamma EQM already represents adequately. The
threshold was not formally optimised nor diagnosed with the
standard residual-life plots; its empirical effect on the
headline metrics is reported in
Section~\ref{sec:results-sensitivity} and the limitation is
acknowledged in the chapter close.

The 80th-percentile choice carries a hidden assumption: every
pixel and dekad combination has enough wet days for stable GPD
fitting. Whether this holds depends on the local wet-day
fraction. To obtain about 30 exceedances at the 80th percentile,
a pixel needs about 150 wet days in the pooled 250-day sample
(since $30 / 0.20 = 150$). Wet-day fractions in Indonesia vary
dramatically with season and region: wet-season Java sees
approximately 70 to 80~percent wet days, giving 175 to 200 wet
days and 35 to 40 exceedances, a comfortable margin. Dry-season
eastern Nusa Tenggara sees approximately 10 to 20~percent wet
days, giving 25 to 50 wet days and only 5 to 10 exceedances,
well below the practical floor for a stable MLE.

In the latter case the GPD fit returns a degenerate default
(shape~$= 0$, location~$= 0$, scale~$= 1$) that produces the
trivial mapping $y = x$, effectively disabling the tail
adjustment for that pixel and leaving the Gamma EQM value in
place. The fallback is defensible in genuinely dry pixels, where
the few wet days are climatologically light-rain events and the
Gamma EQM result is the right answer by default; the fallback
is less defensible in pixels with 25 to 50 wet days, where heavy
events do exist but cannot be fit reliably. The fraction of
pixel-dekads in which the GPD adjustment is active versus
fallback was not reported as a spatial diagnostic in the present
study; producing such a map would be a useful supplementary
artefact for future work.

\medskip\noindent\textbf{Why subsample-averaged MLE rather than a single fit.}
A single GPD maximum-likelihood fit to approximately 50
exceedances has wide confidence intervals on both $\xi$ and
$\sigma$. The rule of thumb in \citep{ref:coles2001} is that GPD
parameter uncertainty becomes operationally tolerable around 100
or more exceedances; with 50 points the single-fit estimate is
noisy and sensitive to a small number of high-leverage values.
A subsample-averaging step reduces estimator variance: the
exceedances are partitioned into five folds of about 10 points
each, the GPD is fitted by maximum likelihood on four folds at a
time (40 points), the parameters $(\xi, \sigma)$ are recorded,
the partition rotates across all five subsamples, and the final
estimate is the mean of the five parameter pairs. The procedure
resembles the subsample step of a bootstrap and reduces variance
arising from a single high-leverage exceedance, but it is not
itself a cross-validation in the strict statistical sense (no
held-out evaluation is performed) and it does not eliminate the
small-sample bias of the GPD shape MLE \citep{ref:hoskingwallis1987}.
The trade is a deliberate one: at 40 to 60 exceedances the
shape estimator is biased but the bias is comparable in
magnitude across folds, so averaging suppresses noise without
making the bias worse.

\medskip\noindent\textbf{Implementation notes.}
The threshold is computed per pixel on the satellite distribution,
not domain-wide and not on the reference. Each pixel uses its
own 80th percentile of its own wet sample. A pixel in Java
might have a threshold around 25~mm/day; a pixel in eastern
Nusa Tenggara might have a threshold of about 8~mm/day. This
per-pixel behaviour is the correct response to the
climatological heterogeneity of the domain. Below the threshold
the Gamma EQM result is used; above it, the GPD-adjusted
quantile is used. There is no smoothing across the join: the
two curves are constrained to match at the threshold percentile,
so the discontinuity is small in practice but is not exactly
zero. The 99.9th-percentile cap is also computed per pixel and
varies spatially: over wet-season Java the cap is on the order
of 200 to 400~mm/day, while over dry-season eastern Nusa
Tenggara it may be 50 to 80~mm/day.


\subsection{Spatial Refinement of Extremes: Convolutional Neural Network}\label{sec:methods-cnn}

The first three stages each operate one pixel at a time. After
the GPD step the corrected field has the right marginal
distribution at every pixel, with the mean from LS, the bulk
shape from Gamma EQM, and the tail from GPD. What it still lacks
is spatial coherence among the extreme pixels. A 100~mm event in
a cluster of adjacent pixels reflects organised convection and
is climatologically plausible; a 100~mm event in a single pixel
surrounded by 5~mm pixels is more likely a satellite retrieval
artefact. Per-pixel statistical correction cannot tell the two
apart, because each pixel is corrected independently.

The fourth stage is a small CNN
that refines extreme-pixel values using local spatial context.
The architecture, implemented in \texttt{src/deep\_learning.py},
is a sequential stack of two Conv2D layers (each followed by
MaxPool2D and Dropout), a Flatten layer, a 128-unit Dense
bottleneck (followed by a third Dropout), and a final Dense
output layer that reconstructs the full grid
\citep{ref:pan2019, ref:banomedina2020}. The output layer uses a
\textit{linear} activation rather than ReLU or sigmoid;
precipitation is non-negative, but a ReLU at the output would
clamp at zero and remove the gradient signal that the mean
squared error loss needs to refine values near the threshold.
The non-negativity constraint is enforced by clipping the
blended output at zero in post-processing; the
99.9th-percentile cap from Section~\ref{sec:methods-gpd} bounds
the upper tail separately and is not re-applied after blending.
Dropout in three places (after each of the two MaxPool layers
and after the Dense bottleneck) provides aggressive
regularisation. The final Dense layer reconstructs the full
grid and dominates the parameter count, so the network is large
relative to the per-dekad training sample; the risk of
overfitting is contained at inference rather than by the
architecture, through applying the refinement only to extreme
pixels, capping its contribution through the blend weight, and
disabling it where station density is low, as discussed below. A
fully-convolutional reconstruction head, prototyped on the Bali
subdomain, matches the dense model's station-validation skill
while cutting the Stage-4 parameter count by roughly $900$ times
at the full Indonesia resolution; the dense head is retained here
so that the released models reproduce the numbers reported in this
thesis, with the fully-convolutional variant available through the
\texttt{architecture} option as the recommended extension for
full-resolution deployment.

The model is trained one per dekadal window on the LSEQM-corrected
field as input and the CPC-UNI reference as target. Loss is
mean squared error over the full field; the extreme-pixel
restriction is applied only at the blending step, not during
training. The Adam
optimizer \citep{ref:kingma2014} drives the parameter update,
and an internal $80/20$ train-validation split with early
stopping (patience of five epochs) controls the number of
training epochs. Thirty-six independent models are produced
and saved for inference. The per-dekad training sample is
approximately 200 days for training and 50 days for validation
after the split. The split is implemented as a random
within-dekad partition rather than a strict temporal hold-out,
so the train and validation samples are drawn from the same
2001 to 2025 range; this is a deliberate compromise that
preserves sample size at the cost of not testing
out-of-distribution generalisation strictly. The batch size
is auto-shrunk at training time if the configured value would
require more than approximately 2~GB per batch on the
$171 \times 461$ Indonesian grid, which it typically does at
default settings.

Three properties of the design constrain the role of the CNN.
First, it operates only above the GPD threshold; sub-threshold
pixels pass through the LSEQM result unchanged. Second, its
output is blended with the LSEQM result rather than replacing
it, under a default weight $\alpha = 0.70$ for LSEQM and
$1 - \alpha = 0.30$ for CNN. Third, the blending weight is
modulated by a station-density confidence mask described in
Section~\ref{sec:methods-blending}, so that the CNN contribution
is suppressed in regions where the CPC reference itself is
weakly constrained by the contributing gauge network. These
three constraints together mean the CNN refines rather than
replaces the statistical baseline: it cannot undo the work of
the earlier stages, and the framework reduces cleanly to LSEQM
when the blend weight is set to $\alpha = 1$, which is the
ablation row reported in the sensitivity analysis of
Section~\ref{sec:results-sensitivity}. The blend weight
$\alpha = 0.70$ and the dependence on station density are
design choices that were not formally optimized; their
sensitivity is reported in
Section~\ref{sec:results-sensitivity}.

Three implementation quirks are worth flagging for any reader
who plans to extend or reuse the CNN component. First, the
choice of mean squared error rather than mean absolute error
penalises large deviations more heavily, which is the right
loss when the goal is to fit the upper tail of the
precipitation distribution; a defensible alternative would be
a quantile loss focused on the same tail, but this was not
tried. Second, the saved model expects a fixed input shape
matching the $171 \times 461$ Indonesian grid; transferring
the framework to a different domain or grid resolution
requires retraining rather than direct inference. Third, a
spatial diagnostic of per-dekad CNN validation loss across the
36 dekadal windows would surface seasons in which the CNN
contribution is weaker because the training extreme-event
sample is thinner (typically dry-season dekads with few
heavy-rain events); such a diagnostic is not formally produced
in the present study and is flagged in
Section~\ref{sec:methods-cnn} as a useful supplementary
artefact for future work.

Figure~\ref{fig:cnn-arch} shows the CNN architecture in detail
alongside the station-density-aware blending mechanism that
Section~\ref{sec:methods-blending} formalises.

% Fig 3.3: dedicated page, image rotated 90 deg to occupy the page height
% as its new width. The original PNG is landscape-oriented; rotating it
% lets the reader see the full diagram at print scale without reducing it
% to thumbnail size on the body text width.
\begin{figure}[p]
\centering
\rotatebox{90}{%
  \includegraphics[width=\figThreeThreeWidth\textheight]{fig_thesis_03_cnn_workflow.png}%
}
\caption[Stage~4 architecture: CNN forward pass, alpha-blending, and station-density confidence mask]{Stage~4 architecture in detail: the Convolutional Neural Network forward pass, the alpha-blending step that combines the CNN refinement with the LSEQM intermediate, and the station-density confidence-mask construction.\label{fig:cnn-arch}}
\end{figure}
\clearpage

Figure~\ref{fig:stage-demo} shows the effect of the four stages
on an illustrative 60-day sample, arranged in a $5 \times 2$
grid read from top to bottom. The reference shown in the figure
is the CPC-UNI calibration target the pipeline acts on; the
independent out-of-sample validation against 172 BMKG stations
is reported in Chapter~\ref{chap:results}. Panels (a, b) show the
raw input and the LS mean match; panels (c, d) show the EQM
quantile-mapping mechanism and its effect on the daily series;
panels (e, f) show the GPD adjustment on the series and on the
wet-day tail; panels (g, h) show the CNN refinement and the final
corrected output; and panels (i, j) summarise the cross-stage
distributions and the headline metrics. The data are synthetic for
illustration; the empirical numbers for the BMKG validation
appear in Chapter~\ref{chap:results}.

Read together, the panels make the division of labour between the
stages concrete. The LS step in panel (b) lifts the whole series
to the gauge mean but leaves the day-to-day spread too narrow, so
the heaviest days stay under-corrected. The EQM step in panels
(c, d) widens the wet-day distribution to match the gauge, yet a
Gamma fit alone still clips the rarest events. The GPD step in
panels (e, f) restores those upper-tail days without disturbing the
bulk of the distribution, and the CNN refinement in panel (g) moves
only the extreme pixels toward the reference. The cumulative effect
in panel (h) is a series whose mean, spread, and tail all track
CPC-UNI; the empirical distributions in panel (i) and the metric
summary in panel (j) confirm that relative bias, the
standard-deviation ratio, and the 99th-percentile ratio each move
toward their targets across the four stages.

\begin{figure}[H]
\centering
\includegraphics[width=\figThreeFourWidth\textwidth]{fig_thesis_03_stage_demonstration.png}
\caption[Illustrative effect of the LSEQM+DL pipeline stages applied to a synthetic 60-day daily-precipitation sample]{Illustrative effect of the LSEQM+DL pipeline stages applied to a synthetic 60-day daily-precipitation sample.\label{fig:stage-demo}}
\end{figure}


% ------------------------------------------------------------
\section{Station-Density-Aware Blending}\label{sec:methods-blending}
% ------------------------------------------------------------

The CNN refinement in Stage~4 is trained against CPC-UNI, and
the quality of CPC-UNI itself depends on the local density of
contributing gauges. Over Java and southern Sumatra the
contributing network is dense and CPC-UNI is a well-constrained
reference. Over Papua, Maluku, and eastern Nusa Tenggara the
network is sparse and CPC-UNI is essentially a smooth-field
reconstruction from few stations. Training a CNN to refine
extremes against such a reference, and then applying the CNN at
its nominal weight everywhere, would let reference uncertainty
in sparse regions propagate into the corrected product. The
blending stage addresses this by making the CNN-versus-LSEQM
weight depend on local gauge density rather than on a constant.
This is a design choice introduced to make the framework
deployable beyond Java without per-region hand-tuning, and is
not claimed as a key contribution of the present work.

Let $C(x, y) \in [0, 1]$ denote a station-density confidence
field constructed once for the domain and cached. The final
corrected value at every pixel above the GPD threshold is the
density-weighted blend of the LSEQM result and the CNN refinement:
\begin{equation}
P_{\text{final}}(x, y, t) =
\alpha_{\text{eff}}(x, y) \cdot P_{\text{LSEQM}}(x, y, t)
+ (1 - \alpha_{\text{eff}}(x, y)) \cdot P_{\text{CNN}}(x, y, t),
\label{eq:blend}
\end{equation}
where the effective blend weight $\alpha_{\text{eff}}$ modulates
the nominal value $\alpha = 0.70$ by the confidence field:
\begin{equation}
\alpha_{\text{eff}}(x, y) = 1 - C(x, y) \cdot (1 - \alpha).
\label{eq:alpha-eff}
\end{equation}
Equation~\ref{eq:alpha-eff} has the following behaviour across
the confidence range:
\begin{itemize}
    \item $C = 1$ gives $\alpha_{\text{eff}} = 0.70$, the
    nominal $30\%$ CNN contribution where the reference is
    well-constrained by the contributing gauge network.
    \item $C = 0.5$ gives $\alpha_{\text{eff}} = 0.85$, a
    $15\%$ CNN contribution at intermediate gauge density.
    \item $C = 0$ gives $\alpha_{\text{eff}} = 1.00$, pure
    LSEQM with the CNN fully suppressed where the reference is
    weak.
\end{itemize}
The relationship between confidence and effective weight is
linear and continuous across this range. The blend is applied
only to pixels above the GPD threshold from
Section~\ref{sec:methods-gpd}; sub-threshold pixels pass
through the LSEQM result unchanged.

The confidence field $C(x, y)$ is built in three steps. First,
the BMKG gauge-station locations of
Section~\ref{sec:methods-data} are counted within each
$0.5^{\circ}$ CPC cell as a proxy for how densely the reference
is constrained there. The same station archive supplies the
out-of-sample validation points, so the deep-learning
refinement is most active in the cells where validation stations
are densest; because the refinement is capped and acts only on
extreme pixels, gauge values are not transferred into the
validation comparison, but this coupling between the density
gate and the validation network is noted here.
Second, the count field is smoothed with
a Gaussian filter of standard deviation $\sigma = 1$ cell, which
corresponds to an e-folding distance of approximately 55~km at
0.5$^{\circ}$ resolution. The smoothing scale is matched to the
CPC-UNI interpolation scale so that the confidence field is
defined at the spatial frequency at which the reference itself
is resolved. Third, the smoothed count is divided by a
saturation parameter and clipped to the unit interval:
\begin{equation}
C(x, y) = \min\!\left(
\frac{(\text{count} \ast G_{\sigma=1})(x, y)}{c_{\text{sat}}},
\; 1
\right),
\label{eq:confidence}
\end{equation}
where $G_{\sigma=1}$ is the Gaussian smoothing kernel and
$c_{\text{sat}} = 2$ is the saturation parameter used in the
present study. The confidence field is then bilinearly
interpolated from the 0.5$^{\circ}$ CPC grid to the
0.1$^{\circ}$ IMERG grid, matching the BCSD strategy applied
elsewhere in the pipeline. The field is computed once, cached as
a NetCDF artefact at \texttt{data/output/station\_density/}, and
reused for every dekadal window and every method.

Both the nominal blend weight $\alpha = 0.70$ and the saturation
parameter $c_{\text{sat}} = 2$ are sensitivity parameters that
were not formally optimized. An earlier configuration of the
pipeline used $c_{\text{sat}} = 3$; qualitative inspection of
the resulting CNN-active fraction across the domain showed that
the higher value suppressed CNN refinement almost everywhere
outside Java and removed the contribution the CNN was designed
to provide, and the default was therefore reduced to 2 for the
present study. The other endpoint is similarly avoided: a lower
parameter (such as $c_{\text{sat}} = 1$) would extend CNN
influence into regions where the CPC reference is too sparse to
support training, which is the failure mode the
station-density-aware blending design is specifically intended
to avoid. The chosen value of $c_{\text{sat}} = 2$ sits at a
defensible empirical balance for Indonesia, and
Section~\ref{sec:results-sensitivity} reports the metric
response to sweeps across both parameters.


% ------------------------------------------------------------
\section{Validation Framework}\label{sec:methods-validation}
% ------------------------------------------------------------

\subsection{Two-Tier Validation}

The corrected satellite product is evaluated at two tiers. The
first tier compares the corrected satellite against CPC-UNI on
the native CPC 0.5$^{\circ}$ grid; this is an in-sample comparison
because CPC-UNI is the correction reference, and it tests
whether the pipeline successfully matches the training target.
The second tier compares the corrected satellite against the 172
BMKG stations on a per-station basis; this is an out-of-sample
comparison because no BMKG station enters the correction step
directly. All headline skill numbers in Chapter~\ref{chap:results}
are computed at the second tier, and the in-sample comparison
is used only as a sanity check that the pipeline did not fail to
fit its training target.

The independence framing in Section~\ref{sec:methods-data} carries
through here: some BMKG stations are present in the CPC-UNI
contributing list, so the two networks are not statistically
independent in a strict sense, but the second tier compares
satellite values extracted at the BMKG station pixel against the
station's own daily record, bypassing the CPC interpolation
entirely. The metric is therefore answering whether the
corrected satellite reproduces the original point observations,
not whether it reproduces the gridded reference.

\subsection{Dekadal Pooling}

The 2001 to 2025 study window is partitioned into 36 dekadal
windows per calendar year: three windows per month, with the
first comprising days 1 to 10, the second days 11 to 20, and
the third days 21 to the end of the month. The 25 years are
pooled within each dekadal window to give approximately 250
paired daily observations per pixel and per window. The
pooling serves two purposes. First, it gives the
distribution-fitting steps in Stages~2 and 3 a sample size
above the operational floor for stable Gamma and GPD fits
\citep{ref:wilks2011, ref:hoskingwallis1997}. Second, it produces
verification statistics whose noise floor is small enough for
between-stage differences to be interpretable: with $n \approx
250$ paired daily observations the Fisher $z$-transform standard
error on Pearson correlation is approximately $\pm 0.06$, which
is well below the inter-stage shifts the pipeline is designed to
produce on the amplitude-related metrics.

Within a dekadal window the partitioning trades inter-annual
variability against intra-dekad seasonal variability. The
implicit assumption is that the bias structure within a fixed
dekadal window is statistically stationary across the 25 years.
This assumption is examined directly in
Section~\ref{sec:results-temporal} (temporal stability across
dekads), which reports the inter-dekad variability of the
headline metrics.

\subsection{Station-Level Extraction}

For each BMKG station, the corrected satellite series is
extracted at the station's reported latitude and longitude on
the 0.1$^{\circ}$ IMERG grid; no spatial averaging is applied.
The satellite day~$t$ is paired with the station day~$t$ in
calendar terms, with no temporal offset beyond the UTC
aggregation already applied to the native half-hourly IMERG
inputs. The station daily values are filtered to remove the
sentinel-coded missing days (Section~\ref{sec:methods-data})
before pairing. The resulting paired series feeds every
out-of-sample metric reported in Chapter~\ref{chap:results}.

\subsection{Reporting Conventions}

Headline numbers in Chapter~\ref{chap:results} are pooled
across the 172 stations and the 36 dekadal windows unless stated
otherwise. Per-region and per-pillar breakdowns are reported in
the appropriate subsections of Chapter~\ref{chap:results} where
the spatial or pillar-level pattern is the object of analysis.
The compact 10-metric subset used for the headline table in
Section~\ref{sec:results-headline} is described in
Section~\ref{sec:methods-metrics}; the full 31-metric catalogue
underlies the Composite Quality Index spatial maps in
Section~\ref{sec:results-cqi}.


% ------------------------------------------------------------
\section{Verification Metric Catalogue}\label{sec:methods-metrics}
% ------------------------------------------------------------

The framework computes 31 pixel-level verification metrics
grouped into four pillars: \textit{value adjustment},
\textit{distribution alignment}, \textit{extreme preservation},
and \textit{event detection} \citep{ref:wilks2011, ref:ebert2007}.
A compact 10-metric subset summarises the headline skill in
Section~\ref{sec:results-headline}; the full catalogue feeds the
Composite Quality Index spatial maps in
Section~\ref{sec:results-cqi}. Pearson correlation between paired
daily values is reported separately as a temporal-skill
diagnostic, not as part of the four pillars; the reason for this
separation becomes the centerpiece of Chapter~\ref{chap:results}
and Chapter~\ref{chap:discussion}.

\subsection{Value Adjustment}

Two metrics test whether the corrected satellite reproduces the
gauge first and second moments. The relative bias (RB) compares the
sample means:
\begin{equation}
\text{RB} = \frac{\overline{P}_{\text{sat}} - \overline{P}_{\text{obs}}}{\overline{P}_{\text{obs}}},
\label{eq:rb}
\end{equation}
with target value zero. The standard deviation ratio (SDR)
compares the sample spreads:
\begin{equation}
\text{SDR} = \frac{\sigma_{\text{sat}}}{\sigma_{\text{obs}}},
\label{eq:sdr}
\end{equation}
with target value unity. RB is the metric the LS stage is
designed to move; SDR is the metric the EQM and CNN stages are
expected to move toward unity.

\subsection{Distribution Alignment}

The distribution-alignment pillar tests whether the corrected
satellite series and the gauge series come from the same
distribution. The Kolmogorov-Smirnov (KS) test
\citep{ref:wilks2011} compares the empirical cumulative
distribution functions $\hat{F}_{\text{sat}}$ and
$\hat{F}_{\text{obs}}$ through their maximum vertical separation:
\begin{equation}
D_{\text{KS}} = \sup_{x}
\left| \hat{F}_{\text{sat}}(x) - \hat{F}_{\text{obs}}(x) \right|,
\label{eq:ks}
\end{equation}
and the reported value is the associated $p$-value, with the
convention that a $p$-value above $0.05$ indicates distributions
that are not statistically distinguishable at the $95\%$
confidence level. The wet-day frequency ratio (WF) tests whether
the corrected series reproduces the gauge fraction of rainy days,
\begin{equation}
\text{WF} = \frac{f_{\text{wet,sat}}}{f_{\text{wet,obs}}},
\label{eq:wetfreq}
\end{equation}
with target unity, where $f_{\text{wet}}$ is the fraction of days
exceeding the $1$~mm/day wet-day threshold. The wet-day intensity
ratio (WI) tests the mean rainfall on wet days,
\begin{equation}
\text{WI} = \frac{\overline{I}_{\text{wet,sat}}}
                 {\overline{I}_{\text{wet,obs}}},
\label{eq:wetint}
\end{equation}
again with target unity, where $\overline{I}_{\text{wet}}$ is the
mean precipitation over wet days only. Six percentile ratios at
the 10th, 25th, 50th, 75th, 90th, and 95th wet-day percentiles
complete the distributional comparison; each is the ratio of the
corresponding quantile,
\begin{equation}
R_q = \frac{Q_q^{\text{sat}}}{Q_q^{\text{obs}}},
\label{eq:pctlratio}
\end{equation}
with target unity at every $q$.

\subsection{Extreme Preservation}

The extreme-preservation pillar focuses on the upper tail of the
distribution, which is the segment that drives flood and drought
thresholds in operational use. The $Q_{95}$ and $Q_{99}$ ratios
are the percentile ratio of Equation~\eqref{eq:pctlratio}
evaluated at $q = 95$ and $q = 99$, with target unity. In
addition, multi-threshold verification at seven thresholds spanning the
operational range (1, 5, 10, 20, 50, 100, and 150 mm/day) reports
the categorical skill of the corrected series, following the
multi-threshold verification approach of the World Meteorological
Organization (WMO) standard TD-1485 \citep{ref:wmo1485}. At each threshold, hits $H$, false alarms
$F$, misses $M$, and correct negatives $Z$ are counted from the
$2 \times 2$ contingency table, from which the threshold-resolved
probability of detection (POD), critical success index (CSI), and
false alarm ratio (FAR) follow by Equations~\eqref{eq:pod}--\eqref{eq:far}
applied at that threshold. The equitable threat score (ETS),
which removes the hits expected by chance, completes the
threshold-resolved set:
\begin{equation}
\text{ETS} = \frac{H - H_{\text{rand}}}{H + M + F - H_{\text{rand}}},
\qquad
H_{\text{rand}} = \frac{(H + M)(H + F)}{N},
\label{eq:ets}
\end{equation}
where $N = H + M + F + Z$ is the total number of day-pixel pairs
at the threshold and $H_{\text{rand}}$ is the expected hit count
under random forecasting. ETS has target unity and a no-skill
value of zero.

\subsection{Event Detection}

The event-detection pillar reports the contingency-table-based
categorical metrics aggregated across all wet days, that is,
without an explicit threshold dependence:
\begin{align}
\text{POD} &= \frac{H}{H + M}, \label{eq:pod} \\
\text{CSI} &= \frac{H}{H + M + F}, \label{eq:csi} \\
\text{FAR} &= \frac{F}{H + F}, \label{eq:far}
\end{align}
where $H$ is the number of hits, $M$ misses, and $F$ false
alarms, all counted at the wet-day threshold of the reference.
Target values are unity for POD and CSI, zero for FAR.
Threshold-resolved versions of these metrics, reported in
Section~\ref{sec:results-threshold}, allow the trade-off between
light-rain and heavy-rain skill to be quantified separately.

\subsection{Composite Quality Index}

The 31-metric catalogue is aggregated into a single per-pixel
Composite Quality Index (CQI) for spatial mapping and qualitative
comparison across regions and methods. The CQI used in this
study is the weighted aggregate of three sub-scores defined in
the source code (\texttt{src/qa\_framework.py}): a Basic
Statistical Quality sub-score (BSQS), a Distribution Quality
sub-score (DQS), and a Temporal Quality sub-score (TQS). Each
sub-score is a weighted sum of normalised component scores
(each mapped to the $[0, 1]$ range, with $1$ best):
\begin{align}
\text{BSQS} &= 0.30\,S_{\text{RB}} + 0.30\,S_{\text{RMSE}}
             + 0.40\,S_{\text{NSE}}, \label{eq:bsqs} \\
\text{DQS}  &= 0.45\,S_{\text{extreme}} + 0.25\,S_{\text{general}}
             + 0.15\,S_{\text{var}} + 0.15\,S_{\text{KS}},
             \label{eq:dqs} \\
\text{TQS}  &= 0.40\,S_{\text{CSI}} + 0.30\,S_{\text{event}}
             + 0.30\,S_{\text{spell}}, \label{eq:tqs}
\end{align}
where $S_{\text{extreme}}$ aggregates the $Q_{90}/Q_{95}/Q_{99}$
percentile-ratio scores (with $Q_{99}$ weighted most heavily),
$S_{\text{general}}$ the $Q_{25}/Q_{50}/Q_{75}$ scores,
$S_{\text{var}}$ the standard-deviation-ratio score,
$S_{\text{KS}}$ the KS-test score, and $S_{\text{event}}$ is the
combination $0.6\,\text{POD} + 0.4\,(1 - \text{FAR})$. The three
sub-scores combine into the CQI with the component weights set in
\texttt{config.yml}:
\begin{equation}
\text{CQI} = 0.40\,\text{BSQS} + 0.30\,\text{DQS}
           + 0.30\,\text{TQS}.
\label{eq:cqi}
\end{equation}
The CQI construction follows the parallel manuscript
\citep{ref:istanto2026mdpi}\MDPIstatus and is reused here without
modification, with the component weights taken from the
configuration file that the pipeline executes
($0.40/0.30/0.30$) rather than the bare function-signature
default.

The three-sub-score CQI and the four-pillar verification
catalogue serve different purposes and group metrics
differently. The four-pillar catalogue is a reporting structure:
it groups metrics by the \textit{bias dimension} each metric
diagnoses (value adjustment, distribution alignment, extreme
preservation, event detection) and supports the stage-by-stage
attribution argument in Chapter~\ref{chap:results}. The CQI is a
\textit{spatial aggregation} object: its three sub-scores group
metrics by the type of statistical question each answers (basic
moments, distributional fit, temporal categorical skill) and
support the spatial mapping in
Section~\ref{sec:results-cqi}. The two groupings overlap, for
example $Q_{99}$ appears in the extreme-preservation pillar and
in the DQS sub-score, but they are not identical and need not
be: the headline numbers reported in
Section~\ref{sec:results-headline} use the four-pillar grouping,
while the spatial maps in Section~\ref{sec:results-cqi} use the
three-sub-score CQI.

\subsection{Pearson Correlation as a Separate Diagnostic}

The per-pixel Pearson correlation between the paired daily
satellite and station observations is reported separately from
the four pillars. For paired daily values $(x_t, y_t)$ at the
same pixel,
\begin{equation}
r = \frac{\sum_t (x_t - \bar{x})(y_t - \bar{y})}
         {\sqrt{\sum_t (x_t - \bar{x})^2}\,
          \sqrt{\sum_t (y_t - \bar{y})^2}},
\label{eq:pearson}
\end{equation}
with target unity. Two further temporal-skill diagnostics
accompany it: the root mean square error (RMSE),
\begin{equation}
\text{RMSE} = \sqrt{\frac{1}{n}\sum_t (x_t - y_t)^2},
\label{eq:rmse}
\end{equation}
with target zero, and the Nash-Sutcliffe efficiency (NSE)
\citep{ref:nashsutcliffe1970},
\begin{equation}
\text{NSE} = 1 - \frac{\sum_t (y_t - x_t)^2}
                      {\sum_t (y_t - \bar{y})^2},
\label{eq:nse}
\end{equation}
with target unity, where $x_t$ is the corrected satellite value,
$y_t$ the gauge observation, and $\bar{y}$ the gauge mean. The
reason for separating these three from the four pillars is
structural: each measures the day-by-day pairing between two
series and is therefore a temporal-skill diagnostic, whereas the
four pillars measure properties of the marginal distribution at
each pixel without regard to the day-by-day pairing. Because the
LSEQM+DL pipeline is a marginal-distribution correction, it is
expected to move the four pillars and to leave Pearson
correlation bounded by the day-by-day pairing of the raw IMERG
input. The empirical verification of this expectation is the
core finding of Chapter~\ref{chap:results}, and its theoretical
basis is developed in Chapter~\ref{chap:discussion}.

Sampling uncertainty on Pearson correlation is reported using
the Fisher $z$-transform standard error
$\text{SE}(z) = 1/\sqrt{n - 3}$, evaluated at $n \approx 250$ to
give $\pm 0.06$. This uncertainty band frames the inter-stage
shifts in Pearson correlation reported in
Section~\ref{sec:results-correlation}.


\subsection{Window-Offset Diagnostic}\label{sec:methods-window}

The Pearson diagnostic of Equation~\ref{eq:pearson} is bounded by
the day-by-day pairing of the satellite and gauge series, which
depends on the 24-hour window used to accumulate each daily total.
The daily IMERG-L archive accumulates over the UTC day, whereas
the BMKG and CPC-UNI gauges accumulate over the local day ending
at the morning observation. To quantify the contribution of this
window definition to the correlation, a window-offset diagnostic
re-aggregates the half-hourly IMERG-L source \citep{ref:imerghhl07}
into 24-hour windows starting at each integer-hour offset
$h \in [-48, +48]$ from $00{:}00$~UTC. For target date~$d$ the
satellite daily total at offset~$h$ is
\begin{equation}
S_d(h) = \tfrac{1}{2}\sum_{k=0}^{47}
         P\!\left(d_0 + h + \tfrac{k}{2}\,\text{h}\right),
\label{eq:window}
\end{equation}
where $P(\cdot)$ is the half-hourly rain rate in mm/h, the factor
$\tfrac{1}{2}$ converts each 30-minute rate to a depth, and $d_0$
is $00{:}00$~UTC on date~$d$. The offset that maximises the
correlation of Equation~\ref{eq:pearson} between $S_d(h)$ and the
BMKG daily total, $h^{\star} = \arg\max_h r(h)$, identifies the
satellite window best aligned with the gauge reporting convention.
The procedure follows the gauge-timing inference of
\citet{ref:beck2019}. The gauge observation time is not altered; only
the satellite aggregation window is shifted. Because the BMKG total
labelled with a given date covers the 24-hour period ending at that
morning's observation, it mostly reflects rain that fell on the
previous calendar day, so the matching satellite window begins
roughly a day before the UTC archive does
(Figure~\ref{fig:window-schematic}).

\begin{figure}[H]
\centering
\includegraphics[width=\figThreeWindowWidth\textwidth]{fig_thesis_03_window_schematic.png}
\caption[Why a $-23$ h window aligns the satellite with the BMKG gauge day]{Why a $-23$~h window aligns the satellite with the BMKG gauge day.\label{fig:window-schematic}}
\end{figure}

Because marginal correction preserves the rank order of the daily
series, the correlation is essentially invariant to the LSEQM+DL
stages (Section~\ref{sec:discussion-bound}). The diagnostic is
therefore computed on the raw half-hourly retrieval rather than on
the corrected product, which isolates the window contribution
without re-running the pipeline at each of the $97$ offsets. The
correlation is pooled through additive sufficient statistics, so
that any aggregation, by pooled station-day pairs, by timezone
band, by calendar month, or by three-month running season, follows
without recomputation. Two poolings are reported: the full BMKG
record ($2001$ to $2021$, the span of the half-hourly archive
processed for the window diagnostic), which resolves the
dependence of the recovery on the satellite era, and the GPM era ($2015$ to $2021$),
which gives the seasonal and per-station climatology over the
years in which IMERG resolves the diurnal cycle. Timezone bands
use the operational WIB, WITA, and WIT assignment ($\text{UTC}+7$,
$+8$, $+9$) recorded for each station.

A monthly re-validation aggregates the daily totals at the
UTC-day and the matched windows into monthly sums and compares
them against the BMKG monthly totals. A station-month is retained
only when no more than five daily values are missing, so that each
monthly total rests on near-complete coverage. This tests whether
the window dependence of the daily correlation survives temporal
aggregation.


% ------------------------------------------------------------
\section{Reproducible Implementation Design}\label{sec:methods-reproducibility}
% ------------------------------------------------------------

The first three sections describe the pipeline as an algorithm.
This section describes it as a software artefact. Reproducibility
is treated here as a parallel methodological contribution of the
thesis: every result reported in Chapter~\ref{chap:results} maps
to a specific function in the source codebase, a specific cell
in a numbered notebook, and a specific configuration value in
\texttt{config.yml}. The empirical demonstration of end-to-end
re-runnability is given in
Section~\ref{sec:results-reproducibility}; this section
describes the design that supports it.

\subsection{Module Architecture}

The Python implementation is partitioned into ten modules under
\texttt{src/}, each with a single responsibility:
\texttt{config} (path and parameter resolution),
\texttt{io} (NetCDF and CSV input and output),
\texttt{utility} (land-mask, lat-flip, sentinel handling),
\texttt{bias\_correction} (LS and EQM stages),
\texttt{distribution\_fitting} (Gamma and GPD parameter
estimation via maximum likelihood, with the subsample-averaging
aggregator for small exceedance samples),
\texttt{deep\_learning} (CNN architecture, training loop, and
inference),
\texttt{station\_density} (confidence-mask construction),
\texttt{metrics} (the 31-metric verification catalogue),
\texttt{qa\_framework} (Composite Quality Index assembly), and
\texttt{station\_validation} (per-station extraction and
pairing). Each module is importable and exposes a small number
of named entry points; no inline function definitions appear in
the notebooks.

\subsection{Configuration-Driven Design}

All paths and parameters are declared in a single YAML file at
the project root, \texttt{config.yml}. A loader,
\texttt{initialize\_config()}, resolves placeholders such as
\texttt{\{main\_dir\}}, \texttt{\{input\_dir\}}, and
\texttt{\{output\_dir\}} to absolute paths at run time. The
sensitivity parameters acknowledged earlier in this chapter
($\alpha = 0.70$, the 80th-percentile GPD threshold, and
$c_{\text{sat}} = 2$) are exposed as named entries in
\texttt{config.yml}, so a sensitivity sweep is a configuration
edit followed by a re-run, with no code changes. The same
configuration file controls the dataset paths, the output
directory layout, the metric pillars, the CNN architecture
hyperparameters, and the per-method output filenames.

\subsection{Notebook Workflow}

The pipeline is exposed through six numbered Jupyter notebooks
under \texttt{notebooks/}: 01 (data acquisition), 02 (LSEQM+DL
bias correction), 03 (performance measurement), 04 (QA
framework), 05 (QA visualisation), and 06 (station validation).
Each notebook follows a
consistent style: short Markdown section headers, small code
cells with a single concern per cell, and calls into the
\texttt{src/} modules rather than inline implementations. The
batch-execution cells at the end of each notebook loop over the
36 dekadal windows with \texttt{try-except} guards so that a
failure in one window does not interrupt the others. The same
notebook can be opened in a local Jupyter installation or in
Google Colab; the behaviour is the same because the configuration
file is identical across environments.

\subsection{Documentation Site}

The \texttt{docs/} directory is a Quarto project that serves as
the documentation site for the framework. The site is structured
into sections that mirror this chapter: \textit{get-started},
\textit{methodology}, \textit{implementation}, \textit{technical},
\textit{example (Bali)}, and \textit{FAQ}. Each documentation page
is cross-linked to the specific notebook cell and source-code
function it describes, so a reader following the documentation can
move directly to the runnable artefact and back. The source code,
notebooks, and documentation are released openly on GitHub under
the Mozilla Public License 2.0, and the full Indonesia input and
output data are deposited on Zenodo with a citable DOI
(Appendix~\ref{app:code}); the documentation site is the
recommended entry point for any reader who wants to run or extend
the pipeline.

\subsection{Environment Specification}

The Python environment is declared in \texttt{environment.yml}
with pinned versions for the core scientific packages
(\texttt{numpy}, \texttt{xarray}, \texttt{scipy},
\texttt{scikit-learn}, \texttt{tensorflow}, \texttt{keras},
\texttt{lmoments3}). The environment was tested on a Windows
local installation with a Conda environment named
\texttt{climate}. Hardware-dependent paths in the CNN training
step are auto-detected at runtime.

Together, these five elements (module architecture, configuration
file, numbered notebooks, documentation site, and
pinned environment) form the reproducibility design of the
framework. Section~\ref{sec:results-reproducibility} demonstrates
that an end-to-end re-run from the published source reproduces
the headline numbers in Chapter~\ref{chap:results} to within the
floating-point tolerance of the underlying libraries.


% ------------------------------------------------------------
\section*{Chapter summary}
% ------------------------------------------------------------

This chapter described the LSEQM+DL pipeline in algorithmic
form and as a software artefact. The four-stage pipeline matches
one correction stage to each of the three marginal-distribution
bias dimensions (LS for the mean, EQM with Gamma for the bulk,
GPD for the tail) and adds a spatial-refinement stage (CNN) for
the extreme pixels. A station-density-aware blending step
modulates the CNN contribution by local gauge density, so the
pipeline degrades gracefully to LSEQM in regions where the
reference is weak. The implementation is partitioned into
single-responsibility modules, driven by a single configuration
file, exposed through six numbered notebooks, and documented
through a Quarto site. Validation against 172 out-of-sample BMKG stations
pools across 36 dekadal windows and 25 years to give the sample
size needed for stable distribution fits and small noise floors
on the headline metrics. The 31-metric verification catalogue is
grouped into four pillars and a separately-reported temporal-skill
diagnostic; the four pillars are the metrics the pipeline is
designed to move, and the temporal-skill diagnostic is the metric
the pipeline is not designed to move. The Chapter~\ref{chap:results}
results are organised around this division.
